% Generated by docs/generate_user_guide.py. Do not edit directly.
% Source: utilities/descriptions.yaml: approaches; models registry

\begin{longtable}{@{}>{\raggedright\arraybackslash}p{0.18\textwidth}>{\raggedright\arraybackslash}p{0.46\textwidth}>{\raggedright\arraybackslash}p{0.28\textwidth}@{}}
\toprule
\textbf{Approach} & \textbf{Description} & \textbf{Available models} \\
\midrule
\endfirsthead
\toprule
\textbf{Approach} & \textbf{Description} & \textbf{Available models} \\
\midrule
\endhead
Approach 1: Direct Correlation & Approach 1 fits a direct relationship between matched in vitro and in vivo time points after preprocessing. The fitted model is then used as a time-scaling function that maps an in vitro time to a predicted in vivo time. & Double Exponential, Exponential, Linear, Linear (zero intercept), Power \\
Approach 2: Independent Model Fitting & Approach 2 fits the same model family independently to the in vitro and in vivo datasets. With fitting data alone, the report summarizes the two fitted models and their fitted-curve uncertainty. When prediction data are supplied, predictions can be generated by comparing the fitted in vivo and in vitro responses or by comparing fitted characteristic time constants when those values are available. & Double Exponential, Exponential, Mass Loss (Gompertz), Mass Loss (One-Phase), Mass Loss (Surface Erosion), Mass Loss (Two-Phase), Power, Stretched Exponential \\
Approach 3: Direct Mapping & Approach 3 uses direct interpolation-based mappings between datasets rather than fitting a parametric model. It can describe value-ratio behavior over time and time-ratio behavior over value ranges. When prediction data are supplied, those mappings can be used to scale prediction values or prediction times. & No parametric model selection \\
\bottomrule
\end{longtable}
